\documentclass[11pt]{article}
\usepackage[margin=1in]{geometry}
\usepackage{amsmath,amssymb,amsthm,mathtools}
\usepackage[hidelinks]{hyperref}
\usepackage{enumitem}
\newtheorem{theorem}{Theorem}[section]
\newtheorem{proposition}[theorem]{Proposition}
\newtheorem{corollary}[theorem]{Corollary}
\newtheorem{lemma}[theorem]{Lemma}
\theoremstyle{definition}
\newtheorem{definition}[theorem]{Definition}
\newtheorem{example}[theorem]{Example}
\newtheorem{remark}[theorem]{Remark}
\newcommand{\Pow}{\mathcal P}
\newcommand{\Fin}{\mathcal P_{\mathrm{fin}}}
\newcommand{\Ext}{\mathsf{Ext}}
\newcommand{\Sat}{\mathsf{Sat}}
\newcommand{\Safe}{\mathsf{Safe}}
\newcommand{\cl}{\operatorname{cl}}
\newcommand{\BA}{\operatorname{BA}}
\newcommand{\im}{\operatorname{im}}
\newcommand{\id}{\operatorname{id}}
\title{Generated Observations and Conservative Characterization Formation}
\author{}
\date{}
\begin{document}
\maketitle

\begin{abstract}
Adding a logically defined test can preserve observational distinctions while
changing a concept lattice or a relational topology. This note separates
syntax, exact semantics, and incomplete evidence, and gives necessary and
sufficient fixed-domain criteria for preserving the FCA extent closure, the
future-observation quotient, and the object Dowker complex. Their intersection
has a largest extensional completion. On a finite nonempty domain without a
universal test, it consists of the old closed extents except the top extent.
Action pullbacks generate the greatest invariant observational quotient, with
a finite refinement bound but no uniform finite bound on infinite domains.
Finite disjunction closure alone forces a full simplex on active objects.
By contrast, all nonempty proper Boolean-definable tests on at least two
and finitely many observable row types give a sphere; permutation dynamics act on its top
homology by sign. Exact simplicial preservation is distinguished from an
exact homotopy criterion for adjoining one test, proved using acyclicity and
a fundamental-group condition. A harmonic oscillator illustrates useful
operation-generated observations and their measurement assumptions. These
are conditional structural results assembled from established theories,
not a claim of priority, an inference of physical laws from a truth table,
or a general theory of scientific discovery.
\end{abstract}

\section{Scope and the three levels of a test}

The question is how characterizations can generate further characterizations
without confusing a change of description with a change of evidence. A useful
answer must specify what a test means, how it is generated, and what the
resulting representation is required to preserve. This note gives such an
answer for fixed-domain observation systems. Its strongest conclusions are
exact preservation criteria and explicit obstructions, rather than a general
algorithm for scientific invention.

A typed probe is a function $p:H\to V$ with a declared value set $V$.
A generating operation has specified input and output types; for example,
$F:V_1\times\cdots\times V_k\to W$ produces
$F(p_1,\ldots,p_k):H\to W$. An action $u:H\to H$ produces the
pullback probe $p\circ u$. Binary probes have $V=\{0,1\}$ and can be
identified with subsets of $H$. A finite binary incidence table is therefore
one important special case, not the definition of every possible scientific
observation. Real-valued observations can be encoded by threshold predicates,
but exact recovery of a real number generally needs infinitely many thresholds;
a finite threshold table records only its stated resolution.

Fix an ambient many-sorted algebra of semantic probes and allowed typed
operations $\Omega$. The least family containing seed probes $S$ and closed
under those operations is
\[
 \operatorname{Gen}_{\Omega}(S)
 =\bigcap\{A:S\subseteq A,\ A\text{ is }\Omega\text{-closed}\}.
\]
This is the ordinary generated-subalgebra construction. It is extensive,
monotone, and idempotent. For a finitary signature its elements are values
of finite terms. Infinitary operations may need transfinite closure and do
not by themselves supply an effective enumeration procedure.

Three levels remain distinct:
\begin{enumerate}[label=(\roman*)]
\item \emph{Syntax}: terms, names, proofs, provenance, and execution recipes.
\item \emph{Semantics}: their interpreted functions on a declared domain.
\item \emph{Evidence}: what has actually been measured or otherwise warranted
about those interpretations.
\end{enumerate}
Evaluation maps the free typed term algebra to the semantic probe algebra;
its image is the generated family and its kernel identifies extensionally
equal terms. Equality in one model or finite sample is not automatically
provable equivalence in a theory, equality in all admissible models, or
equality of computational cost. For example, $p$ and $p\wedge p$ agree
extensionally while having different syntactic lengths.

Every pointwise-generated probe factors through the joint values of its
input probes. Thus pointwise recombination alone cannot distinguish two
states whose old values agree. Action pullback is different: it uses the
transition map in addition to the current row. The following sections make
that distinction precise and then show that the same point resolution can
support different FCA and topological structures.

\section{Observations and the meanings of preservation}

Let $H$ be a set, $T$ a set of test labels, and
$R\subseteq H\times T$ an incidence relation. Write
\[
 E_t=\{h\in H:hRt\},\qquad \mathcal E=\{E_t:t\in T\}\subseteq\Pow(H).
\]
The family $\mathcal E$ forgets repeated labels with the same extent. Results
about a largest completion below are \emph{extensional}: arbitrarily many
synonymous labels would preclude a largest labeled language.
Define
\begin{align*}
 \Ext(\mathcal E)&=\left\{\bigcap\mathcal F:\mathcal F\subseteq\mathcal E\right\},
 &\bigcap\varnothing&=H,\\
 \cl_{\mathcal E}(A)&=\bigcap\{E\in\mathcal E:A\subseteq E\},\\
 K(\mathcal E)&=\{\varnothing\}\cup
 \{\sigma\in\Fin(H):\exists E\in\mathcal E\ (\sigma\subseteq E)\},\\
 N(T,R)&=\{\varnothing\}\cup
 \{\tau\in\Fin(T):\bigcap_{t\in\tau}E_t\ne\varnothing\}.
\end{align*}
Thus $\Ext(\mathcal E)$ is the FCA extent lattice under inclusion,
$K(\mathcal E)$ is the object Dowker complex, and $N(T,R)$ is its test-side
counterpart. Write $\Delta(D)=\Fin(D)$ for the full abstract simplex.
The empty face is included even when there are no nonempty faces;
its geometric realization is then empty. Ground elements whose singleton is
not a face are not vertices. In the usual FCA notation,
$A'=\{t:A\subseteq E_t\}$ and $B'=\bigcap_{t\in B}E_t$.
The formal concepts are pairs $(A,B)$ with $A'=B$ and $B'=A$;
projection to $A$ identifies their ordered lattice with
$\Ext(\mathcal E)$. Equality of this labeled extent closure is stronger
than an unspecified abstract lattice isomorphism.

Suppose also that a monoid $M$ acts on $H$ on the right. Put
\[
 \tau_m(h)=hm,\qquad U_m(C)=\tau_m^{-1}(C),\qquad
 h\equiv_{\mathcal E}k
 \Longleftrightarrow
 \forall m\in M\ \forall E\in\mathcal E\quad
 (hm\in E\Longleftrightarrow km\in E).
\]
Then $\equiv_{\mathcal E}$ is a right congruence, so
$Q_{\mathcal E}=H/{\equiv_{\mathcal E}}$ carries an induced right $M$-action.
Here $\tau_{mn}=\tau_n\circ\tau_m$ and $U_{mn}=U_m\circ U_n$.
Let $\Sat(\equiv_{\mathcal E})$ denote the subsets of $H$ that are unions of
congruence classes.

An additional family of tests is specified by its extents
$\mathcal G\subseteq\Pow(H)$. The extended incidence has family
$\mathcal E^+=\mathcal E\cup\mathcal G$ on the \emph{same} $H$ and with the
\emph{same} action. Throughout, preservation of the FCA structure means
literal equality of the closure operators, equivalently of their fixed
subsets of $H$. Preservation of the residual structure means equality of the
equivalence relations on $H$, hence of their canonical quotient actions.
Preservation of the object complex means equality of its finite faces.
These assertions do not say that the enlarged labeled vocabulary, its full
implication theory, or its test-side complex is literally unchanged.

\begin{proposition}[Universal behavioral quotient]\label{prop:universal}
For arbitrary $H$ and $M$, the map $q:H\to Q_{\mathcal E}$ is the
coarsest equivariant quotient that preserves the original test readouts.
Precisely, let $Y$ carry a right $M$-action, let $f:H\twoheadrightarrow Y$
be equivariant, and suppose there are subsets $D_t\subseteq Y$ with
$f^{-1}(D_t)=E_t$. There is a unique equivariant surjection
$p:Y\twoheadrightarrow Q_{\mathcal E}$ satisfying $q=p f$.
If $\overline E_t=q(E_t)$, then $p^{-1}(\overline E_t)=D_t$.
For an additional test $C\subseteq H$, a readout
$\overline C\subseteq Q_{\mathcal E}$ with $C=q^{-1}(\overline C)$
exists if and only if $C$ is $\equiv_{\mathcal E}$-saturated.
\end{proposition}
\begin{proof}
If $f(h)=f(k)$ then $f(hm)=f(km)$ for every $m$, so all old tests agree
after every action. Thus $\ker f\subseteq\equiv_{\mathcal E}$, and
$p(f(h))=q(h)$ is well-defined. Surjectivity of $f$ gives uniqueness;
equivariance and preservation of readouts follow by substitution.
The last assertion is exactly the condition for a subset to be an inverse
image under a quotient map.
\end{proof}

\section{Exact criteria for arbitrary simultaneous extensions}

\begin{theorem}[Three extension criteria]\label{thm:criteria}
For arbitrary sets $H,T,M$ and an arbitrary family $\mathcal G$:
\begin{enumerate}[label=(\roman*)]
\item $\Ext(\mathcal E^+)=\Ext(\mathcal E)$ if and only if
      $\mathcal G\subseteq\Ext(\mathcal E)$.
\item $\equiv_{\mathcal E^+}=\equiv_{\mathcal E}$ if and only if
      $\mathcal G\subseteq\Sat(\equiv_{\mathcal E})$.
\item $K(\mathcal E^+)=K(\mathcal E)$ if and only if
      $\Fin(C)\subseteq K(\mathcal E)$ for every $C\in\mathcal G$.
      If each added $C$ is finite (in particular if $H$ is finite), the last condition is equivalent to
      $C=\varnothing$ or $C\subseteq E$ for some $E\in\mathcal E$,
      separately for every $C\in\mathcal G$.
\end{enumerate}
Moreover $\Ext(\mathcal E)\subseteq\Sat(\equiv_{\mathcal E})$.
Consequently FCA preservation already implies residual preservation in this
fixed-state, fixed-action extension problem.
\end{theorem}

\begin{proof}
Each added extent is itself an intersection of a singleton family, so (i)
is necessary. If every added extent was closed, every intersection of old
and new extents is still closed; this proves sufficiency.

Extension can only refine the residual equivalence. Equality therefore
implies that every added $C$ is saturated, by testing the identity of $M$.
Conversely, $h\equiv_{\mathcal E}k$ implies $hm\equiv_{\mathcal E}km$ for
every $m$. Saturation of $C$ then gives $hm\in C\Leftrightarrow km\in C$.
This proves (ii), including all future actions, not just the current rows.

For (iii), the new faces are precisely the finite subsets of added
extents. If $C$ is finite and nonempty, the finite subset $C$
itself must have an old witnessing column. The converse is immediate.
Finally, old columns are saturated by the identity test, and intersections
of saturated sets are saturated.
\end{proof}

\begin{definition}[Common conservative completion]
Set
\[
 \Safe(\mathcal E)=
 \{C\in\Ext(\mathcal E):\Fin(C)\subseteq K(\mathcal E)\}.
\]
\end{definition}

\begin{theorem}[Largest simultaneous completion]\label{thm:maximal}
The family $\Safe(\mathcal E)$ contains $\mathcal E$ and is the unique
largest extensional extension that preserves the three structures in
Theorem~\ref{thm:criteria}. In particular,
\[
 \Ext(\Safe(\mathcal E))=\Ext(\mathcal E),\quad
 K(\Safe(\mathcal E))=K(\mathcal E),\quad
 \equiv_{\Safe(\mathcal E)}=\equiv_{\mathcal E},
 \qquad \Safe(\Safe(\mathcal E))=\Safe(\mathcal E).
\]
Every proper closed extent belongs to $\Safe(\mathcal E)$. The only
possible excluded closed extent is $H$, and
\[
 H\in\Safe(\mathcal E)\quad\Longleftrightarrow\quad
 \Fin(H)\subseteq K(\mathcal E).
\]
If $H$ is finite and nonempty, this is equivalent to $H\in\mathcal E$.
Thus, in the finite nonempty case without a universal test,
\[
 \Safe(\mathcal E)=\Ext(\mathcal E)\setminus\{H\}
 =\left\{\bigcap\mathcal F:\varnothing\ne\mathcal F\subseteq\mathcal E\right\}.
\]
When $H=\varnothing$, $\Safe(\mathcal E)=\{\varnothing\}$.
\end{theorem}

\begin{proof}
The three preservation statements and maximality follow from
Theorem~\ref{thm:criteria}; idempotence follows because the defining closure
family and complex have not changed. If $C$ is a proper closed extent,
write $C=\bigcap\mathcal F$. Since the empty intersection is $H$,
$\mathcal F$ is nonempty. Choose $E\in\mathcal F$; then $C\subseteq E$,
so all finite subsets of $C$ are old faces. The remaining assertions follow
from the definition and the finite criterion in Theorem~\ref{thm:criteria}.
In the finite case without a universal column, an intersection of a
nonempty family of columns cannot be $H$.
\end{proof}

\begin{remark}[What the completion accomplishes]
Nonempty conjunctions introduce names for already available
characterizations while preserving all three reconstructions. The empty
conjunction is different: it introduces a universal test. The theorem
identifies the largest common semantic completion; it does not assert that
this vocabulary is the only scientifically useful one. A deliberately
discriminating experiment or a desired disjunction can properly leave it.
\end{remark}

\begin{proposition}[Complete failure certificates]\label{prop:certificates}
Let $C\subseteq H$. Failure of any one of the three criteria has the
following witness:
\begin{enumerate}[label=(\roman*)]
\item FCA failure: $h\in\cl_{\mathcal E}(C)\setminus C$.
\item Residual failure: $h\equiv_{\mathcal E}k$ with exactly one of
      $h,k$ in $C$.
\item Literal object-complex failure: a finite nonempty
      $\sigma\subseteq C$ contained in no old column.
\end{enumerate}
For finite $H$ these are finite witnesses, relative to the supplied finite
table and action semantics. They need not be certificates about unobserved
states of an unknown system.
\end{proposition}
\begin{proof}
Closure is extensive, so $C$ fails to be closed exactly when the difference
in (i) is nonempty. The other two statements are the negations of
saturation and the finite-face criterion.
\end{proof}

Finiteness of witnesses does not imply effective verification for an
arbitrary presentation of $M$ or an unenumerated test family. With explicit
finite tables for $H$, the tests, and a finite set of action generators,
the refinement construction below computes the residual relation and the
checks are finite. No such algorithmic claim is made for inaccessible
actions merely described abstractly.

\section{Logical definitions and genuinely new observations}

There are two notions of observational indistinguishability that must be
kept separate. Equality of the \emph{current} rows means agreement on the
tests currently available. The residual relation already quantifies over
\emph{all} permitted future actions. Generating executable future tests can
refine the former toward the latter, even though it cannot refine the
already completed residual relation for that same action system.

\subsection{Future experiments generate a stable behavioral quotient}

In deterministic modal syntax, $[a]\varphi$ holds at $h$ exactly when
$\varphi$ holds at $ha$, so its extent is $U_a(\operatorname{Ext}(\varphi))$.
Here $\operatorname{Ext}(\varphi)$ denotes the truth set of the formula.
This is a semantic clause using the supplied dynamics, not an inference
of that dynamics from the current truth table.

Fix a set of action generators $A$ whose generated submonoid is $M$.
For a word $w\in A^*$, use the associated state map $\tau_w$ and inverse
image $U_w$. Set
\[
 \mathcal E^{\le n}=\{U_w(E):E\in\mathcal E,\ |w|\le n\},\qquad
 h\sim_n k\ \Longleftrightarrow\
 \forall C\in\mathcal E^{\le n}\ (h\in C\Longleftrightarrow k\in C).
\]
The original instantaneous row equivalence is $\sim_0$.

\begin{theorem}[Generation by bounded future experiments]\label{thm:future}
For arbitrary $H$,
\[
 \sim_{n+1}=\sim_n\ \cap\!
 \bigcap_{a\in A}(\tau_a\times\tau_a)^{-1}(\sim_n),\qquad
 \bigcap_{n\ge0}\sim_n=\equiv_{\mathcal E}.
\]
The limiting relation is the greatest right $M$-congruence contained in
$\sim_0$. Each pair separated at depth $n$ has a distinguishing original
test preceded by a word of length at most $n$. If it was not separated
at depth $n-1$, its shortest such word has length exactly $n$.

If $|H|=r<\infty$ and $c_0=|H/{\sim_0}|$, the sequence stabilizes after
at most $r-c_0$ strict refinement steps. Once $\sim_{n+1}=\sim_n$, it is
stable forever and already equals $\equiv_{\mathcal E}$. In particular
every behaviorally distinct pair on this finite domain has a separating
future experiment of length at most $r-c_0$.
\end{theorem}

\begin{proof}
A word of length at most $n+1$ either is already of length at most $n$ or
begins with one generator followed by a word of length at most $n$.
This proves the recurrence and, by the definition of $M$, the limiting
equality. The limit is stable under every generator: agreement at every
depth implies agreement after applying a generator at every depth.
Any congruence contained in $\sim_0$ forces agreement after every word,
so is contained in the limit. The experiment witnesses follow directly
from the definitions. On a finite set every strict refinement increases
the number of classes by at least one, which can happen at most
$r-c_0$ times. Equality at one stage is a fixed point of the displayed
recurrence, so a later strict step cannot occur.
\end{proof}

\begin{corollary}[No artificial growth from a Chu-complete language]
If each generator has a Chu translation on the initial tests, then
$\sim_0$ is already a right congruence and
$\sim_n=\sim_0=\equiv_{\mathcal E}$ for every $n$. Independently,
pointwise Boolean definitions in any fixed current test family never
split its current row equivalence. Future pullbacks and pointwise Boolean
definitions are therefore different generation mechanisms.
\end{corollary}
\begin{proof}
Under a Chu translation, agreement on all current tests implies agreement
after applying the generator, because its pullback is again a current
test. Boolean combinations preserve equal truth-value vectors.
\end{proof}

\begin{example}[Infinite future depth]
Let $H=\mathbb N$, let $u(n)=\max(n-1,0)$, and let the only initial test
be $\{0\}$. Its depth-$j$ pullback is $\{0,\ldots,j\}$. At depth $n$,
the classes are $\{0\},\ldots,\{n\}$ and the infinite tail
$\{k:k>n\}$. Thus the limiting residual relation is equality but no
finite stage is stable.
\end{example}

Theorem~\ref{thm:future} is the classical partition-refinement mechanism
of Moore and Myhill--Nerode theory, expressed here as generation of test
extents; see Moore \cite{Moore}. It provides an operationally meaningful
source of new current observations, not a new theorem of automata
minimization. It assumes the specified actions and observations can be
performed and are semantically known. A test splitting an \emph{existing
full residual class} requires a genuinely new observation or an expansion
of the permitted action system, rather than one of its already accounted
future recipes.

\subsection{The quotient must retain the structure of its tests}

Let $\mathcal O=\{U_m(E):m\in M,E\in\mathcal E\}$. Choose one
specified closure doctrine $\mathfrak D$: arbitrary intersections for an
FCA closure system; finite intersections and arbitrary unions for a topology;
Boolean operations; or countable unions and complements for a sigma-algebra.
Let $\mathfrak D(\mathcal O)$ be the least family closed under the chosen
operations. Since inverse images commute with each of them, this is the least
$\mathfrak D$-closed, action-stable family containing $\mathcal E$.
Every operation listed preserves agreement of truth values, so its induced
point equivalence is exactly $\equiv_{\mathcal E}$.

Write $q:H\twoheadrightarrow Q_{\mathcal E}$ and put
\[
 \mathcal L_Q=\{q(C):C\in\mathfrak D(\mathcal O)\}.
\]
All such $C$ are saturated, and therefore
\[
 q^{-1}:\mathcal L_Q\longrightarrow\mathfrak D(\mathcal O)
\]
is a bijection preserving every selected operation; its inverse is direct
image on these saturated subsets. This proves exact preservation of the
\emph{structured} quotient $(Q_{\mathcal E},\mathcal L_Q)$, not just
of its point set. The universal factorization of
Proposition~\ref{prop:universal} also respects this structure whenever the
competing representation carries the corresponding generated probes.

The chosen doctrine matters. On $H=\{0,1\}$, the single extent $\{0\}$
and the two extents $\{0\},\{1\}$ both distinguish the states, but their
FCA lattices have respectively two and four elements. Likewise the
Sierpi\'nski and discrete topologies on two points have the same equality
of observable points and different frames of opens. A bare quotient cannot
replace these structures. This is closely related to complete abstractions
and strong preservation in abstract interpretation \cite{RanzatoTapparo},
and to behavioral semantics in coalgebra \cite{Rutten}.

\subsection{Boolean definitions and disjunction}

\begin{definition}
An \emph{acted definable test} is an extent obtained from finitely many
sets $U_m(E_t)$ by Boolean operations. This includes the constants $H$
and $\varnothing$. A new experiment is \emph{behaviorally discriminating}
relative to $\mathcal E$ if its extent is not
$\equiv_{\mathcal E}$-saturated. Extensional definability is relative to the
specified whole state domain $H$, not merely to a finite observed sample.
\end{definition}

\begin{theorem}[Boolean completion and its cost]\label{thm:boolean}
Suppose $H$ is finite. Let
\[
 \mathcal B=\BA\{U_m(E):m\in M,\ E\in\mathcal E\}.
\]
Then
\[
 \mathcal B=\Sat(\equiv_{\mathcal E}),\qquad
 \equiv_{\mathcal B}=\equiv_{\mathcal E},\qquad
 \Ext(\mathcal B)=\mathcal B\cong\Pow(Q_{\mathcal E}),\qquad
 K(\mathcal B)=\Fin(H).
\]
If $H\ne\varnothing$, the last complex has contractible realization.
Thus a maximal Boolean language can preserve every behavioral distinction
while changing both the FCA structure and the topology. In particular,
logical conservativity does not imply native topological conservativity.
\end{theorem}

\begin{proof}
Each generator is saturated, hence every Boolean combination is saturated.
For $h\in H$ and each $k\not\equiv_{\mathcal E}h$, choose an acted
predicate separating $h$ from $k$, and choose it or its complement to
contain $h$. The finite conjunction of these choices is exactly the
equivalence class of $h$. Every union of classes is a finite Boolean
combination, proving the first equality. Saturated subsets are closed
under all intersections and all action inverse images, proving the next
two assertions. The universal test $H$ lies in $\mathcal B$, and therefore
witnesses every finite face.
\end{proof}

\begin{theorem}[Finite disjunction alone removes Dowker holes]\label{thm:disjunction}
For arbitrary $H$, let $V=\bigcup\mathcal E$ be the active object set and
let
\[
 \mathcal E^\vee=
 \left\{E_1\cup\cdots\cup E_k:k\ge1,\ E_i\in\mathcal E\right\}.
\]
Then $K(\mathcal E^\vee)=\Fin(V)$. If $V\ne\varnothing$, the completed
object complex is contractible, as is its test-side Dowker complex.
No truth constant, infinite disjunction, or finiteness assumption on
$H$ is needed. The instantaneous row equivalence and the full residual
equivalence are both unchanged by this extension.
\end{theorem}

\begin{proof}
Every union is contained in $V$, so every face lies in $\Fin(V)$.
For a finite nonempty $\sigma\subseteq V$, choose for each
$h\in\sigma$ one old column $E_h$ containing it. The finite nonempty
disjunction $\bigcup_{h\in\sigma}E_h$ witnesses $\sigma$ as a face.
The empty face is present by convention. The full complex on nonempty
$V$ is a cone at any chosen vertex, including when $V$ is infinite.
Dowker's homotopy equivalence gives the test-side assertion
\cite{BrunSalbu}. Pointwise unions preserve equality of current rows;
unions of residual-saturated extents are residual-saturated, so
Theorem~\ref{thm:criteria} gives the residual assertion.
\end{proof}

The theorem treats every generated compound predicate as an admissible new
witness element of the incidence relation. It is not a claim that physically
executing the compound test is free, nor that adding it to a geometric cover
preserves good-cover hypotheses. Keeping formulas available for reasoning
without adding them as witness columns is a different construction.

Deleting universal extents afterward does not generally leave a
disjunction-closed language: two proper extents can have universal union.
Moreover this does not avoid the infinite-domain obstruction. For all
singleton columns on $\mathbb N$, every finite disjunction is a proper
finite subset, yet the resulting object complex is already the full
infinite simplex. Nonempty conjunction completion preserves the object
complex exactly; unrestricted finite disjunction completion generally
does not, even without a universal test.

\begin{remark}[Infinite domains]
The equality between finitary Boolean definability and all saturated sets
can fail when $H$ is infinite. For $H=\mathbb N$ with trivial action and
all singleton tests, the residual relation is equality, but their
finitary Boolean algebra consists only of finite and cofinite sets.
Allowing arbitrary unions and intersections yields all saturated sets.
Separately, the finite-column criterion for topology cannot be used for
infinite $C$: for $E_n=\{0,\ldots,n\}$ every finite subset of $\mathbb N$
has an old witness although no $E_n$ contains $\mathbb N$. Adding the
universal test in that example leaves the object complex literally equal.
\end{remark}

For a labeled context, an implication $A\Rightarrow B$, with
$A,B\subseteq T$, holds exactly when
\[
 \bigcap_{t\in A}E_t\subseteq\bigcap_{t\in B}E_t.
\]
Adding test labels while leaving $H$ and the old incidences fixed preserves
all implications expressed solely in the old language, even when the new
extent is \emph{not} closed and the FCA lattice changes. The enlarged
language nevertheless has a different set of implications. This
old-language conservativity is strictly weaker than equality of the extent
closure and should not be substituted for it.

\begin{proposition}[Rows and tests have different evidential effects]
Let $H\subseteq\widehat H$ and extend all old tests to $\widehat H$ while
preserving their values on $H$. Every old-language implication true on
$\widehat H$ was true on $H$. A formerly true implication fails precisely
when an added row satisfies its premises and violates a conclusion.
In contrast, adding columns on fixed $H$ does not change the truth of any
old-language implication.
\end{proposition}
\begin{proof}
Validity is universal quantification over rows. Extending that domain can
add counterexamples but cannot remove any. Adding labels does not change
the domain or interpretation of an old formula.
\end{proof}

This gives a limited but exact co-generation discipline: derived
characterizations generate labeled tests; new rows can refute their
empirical implications; behaviorally discriminating tests split the
residual state space. None of these mechanisms by itself selects a
scientifically worthwhile experiment, supplies its physical realization,
or establishes a law beyond the specified observation domain.

\subsection{Positive formulas preserve an order rather than every concept}

Define $h\preceq k$ when every old predicate true at $h$ is true at $k$.
An added extent preserves this preorder exactly when it is upward closed;
it preserves current row equivalence exactly when it is a union of current
row classes. Closed FCA extents satisfy both conditions, but the converses
fail. On a finite row-type quotient $Q$, all finite positive formulas with
constants define exactly the upsets of this order. Their FCA closure is
\[
 c_{+}(X)=\bigcup_{h\in X}\cl_{\mathcal E}(\{h\}).
\]
Indeed each principal upset is a finite conjunction of predicates excluding
the finitely many types not above it, and every upset is a finite union of
principal upsets. Thus positive completion preserves the old FCA closure
exactly when the old closure operator preserves arbitrary unions, including
the empty union. Monotone
syntax alone is insufficient for FCA preservation. Adding negation instead
produces all subsets of a finite row-type quotient, as described next.

A related punctured construction retains all nonempty proper upsets.
For finite nonempty $Q$, let $\operatorname{Min}(Q)$ be its minimal
elements. Its quotient object complex is
\[
 \bigcup_{m\in\operatorname{Min}(Q)}\Delta(Q\setminus\{m\}).
\]
Every proper upset misses a minimal element, and each displayed complement
is itself an upset. A nonminimal element is a cone vertex, so this complex
is contractible unless $Q$ is an antichain. For an antichain of size $q\ge2$
it is the boundary sphere $S^{q-2}$; for $q=1$ it has empty realization.
Duplicated rows retract onto representatives by the same contiguity argument
used below. This punctured family is not generally disjunction-closed.

\section{Boolean saturation and presentation dependent topology}
\label{sec:saturation}

Here $\sim$ denotes equality of the rows of a specified starting family
$\mathcal P\subseteq\Pow(H)$, and $\pi:H\twoheadrightarrow Q=H/{\sim}$.
It may be the instantaneous family $\mathcal E$, or the full future family
$\mathcal O$; the choice must be stated. Set
$\mathcal B=\BA(\mathcal P)$ and
$\mathcal B^*=\mathcal B\setminus\{\varnothing,H\}$.
All formulas are finitary. Removing constants \emph{by their extents} is
stronger than omitting their symbols: $p\vee\neg p$ is still universal.
Moreover $\mathcal B^*$ is usually not closed under the Boolean operations.

\begin{lemma}[Finite type saturation]\label{lem:finite-types}
If $Q$ is finite, even when $H$ and $\mathcal P$ are infinite,
\[
 \mathcal B=\{\pi^{-1}(S):S\subseteq Q\}.
\]
\end{lemma}
\begin{proof}
Boolean formulas are constant on row classes. For a class $q$, choose
for each $r\ne q$ a literal in an old predicate that is true on $q$ and
false on $r$. The finite conjunction defines exactly $q$; finite unions
of these formulas define every subset of $Q$. The empty conjunction and
disjunction handle the constants and the one-class case.
\end{proof}

\begin{theorem}[Punctured Boolean saturation]\label{thm:sphere}
If $2\le q=|Q|<\infty$, then
\[
 K(\mathcal B^*)=
 \{\sigma\in\Fin(H):\pi(\sigma)\ne Q\},
 \qquad |K(\mathcal B^*)|\simeq S^{q-2}.
\]
There is a strong deformation retraction onto the boundary of the
$(q-1)$-simplex carried by any chosen set of one representative per class.
Arbitrarily many, including infinitely many, duplicate rows cause no change
in this conclusion. The test-side Dowker complex has the same homotopy type.

If $Q$ is infinite, $K(\mathcal B^*)=\Fin(H)$, so its realization is
contractible. No infinitary Boolean operations are needed for this assertion.
If $q=0$ or $q=1$, the punctured complex has empty realization.
\end{theorem}
\begin{proof}
For finite $Q$, Lemma~\ref{lem:finite-types} shows that a nonempty finite
$\sigma$ has a proper Boolean witness exactly when its represented types
omit some member of $Q$. Let $L=\partial\Delta(Q)$, where
$\Delta(Q)=\Fin(Q)$. Projection $r=\pi:K(\mathcal B^*)\to L$ is
simplicial. A chosen representative section $s:Q\to H$ is simplicial
in the opposite direction, and $rs=\id$. For every face $\sigma$,
\[
 \pi\bigl(\sigma\cup sr(\sigma)\bigr)=\pi(\sigma)\ne Q.
\]
Thus $sr$ is contiguous to the identity relative to the representative
subcomplex. The straight-line homotopy on each common simplex gives the
claimed deformation retraction, with the weak CW realization topology.
The boundary of a $(q-1)$-simplex is $S^{q-2}$.

Equivalently, the $q$ full simplices on
$H\setminus\pi^{-1}(\{i\})$ cover $K(\mathcal B^*)$; every proper
subfamily has contractible nonempty intersection and the intersection of
all $q$ is empty. Their nerve is $\partial\Delta(Q)$. On representatives
these coatom tests give the classical contranominal relation $(Q,Q,\ne)$.
The other proper Boolean tests add no object faces.

For infinite $Q$, take a nonempty finite
$\sigma=\{h_1,\ldots,h_n\}$ and choose $y$ with a type outside
$\pi(\sigma)$. For each $h_i$ choose a literal true on $h_i$ and false
on $y$. Their finite disjunction contains $\sigma$ and excludes $y$;
it is a nonempty proper Boolean extent. Hence every finite subset of $H$
is a face. The full complex contracts to any chosen vertex.
If $q=1$, the only Boolean extents are $\varnothing,H$; if $q=0$, then
$H=\varnothing$. In both cases the punctured family has no active tests.
Dowker duality gives the test-side statements \cite{BrunSalbu}.
\end{proof}

For $q=2$ the complex is the disjoint union of full simplices on the two
blocks, so its reduced $H_0$ is $\mathbb Z$. For $q=1$ the abstract
complex is $\{\varnothing\}$, sometimes denoted $S^{-1}$ under augmented
homology conventions, with reduced $H_{-1}=\mathbb Z$. It is not a point.
For $q=0$ no $S^{-2}$ assertion is made. Including truth gives a
contractible full object simplex whenever $H\ne\varnothing$; false is a
ghost test and has no effect on nonempty faces.

\begin{proposition}[Permutation dynamics and singular dynamics]
\label{prop:sign}
Suppose $2\le q<\infty$ and $f:H\to H$ respects row equivalence,
inducing $\sigma:Q\to Q$ with $\pi f=\sigma\pi$.
If $\sigma$ is a permutation, $f$ induces a simplicial endomorphism of
$K(\mathcal B^*)$, and its map on
$\widetilde H_{q-2}(-;\mathbb Z)\cong\mathbb Z$ is multiplication by
$\operatorname{sgn}(\sigma)$. If $\sigma$ is not bijective, $f$ still
induces a simplicial map, but it is nullhomotopic and its reduced-homology
maps are zero.
\end{proposition}
\begin{proof}
A permutation preserves proper type supports, and the canonical projection
intertwines $f$ with the permutation of $\partial\Delta(Q)$. For an
ordering $Q=\{1,\ldots,q\}$, the fundamental cycle is
$z=\partial[1,\ldots,q]$. Its image is
$\partial[\sigma(1),\ldots,\sigma(q)]
 =\operatorname{sgn}(\sigma)z$.
If $\sigma$ is nonbijective, finiteness makes $\sigma(Q)$ a proper
nonempty subset. The whole image of $f$ lies in the full simplex on
$\pi^{-1}(\sigma(Q))$, so the induced map is nullhomotopic.
\end{proof}

No bijectivity of $f$ within duplicate blocks is needed. No equivariant
representative section is asserted: equivariance of the canonical projection
suffices for homology. Over characteristic two the sign is invisible.
Predicate pullback has a separate condition. For a permutation,
$f^{-1}(\pi^{-1}S)=\pi^{-1}(\sigma^{-1}S)$ preserves nonempty proper
tests. For a nonbijective finite quotient self-map, it does not:
$\sigma^{-1}(\sigma(Q))=Q$, while the preimage of a singleton outside
the image is empty. Thus the punctured Boolean family need not be stable
under operational pullback, despite existence of the covariant object map.

\begin{proposition}[The map in a Boolean refinement]\label{prop:refinement}
Let $\mathcal B_t\subseteq\mathcal B_u$ be finite Boolean subalgebras
on the same fixed $H$, with $2\le q_t<q_u$ realized types. The canonical
inclusion $K(\mathcal B_t^*)\hookrightarrow K(\mathcal B_u^*)$ is
nullhomotopic and induces zero on every reduced homology group.
\end{proposition}
\begin{proof}
Choose one object per old type. The old complex retracts onto the boundary
on those representatives. Their $q_t$ new types form a proper subset of
$Q_u$, so that entire representative boundary lies inside a simplex of the
new complex and contracts there. Compose the old retraction homotopy with
this contraction.
\end{proof}

The quotient map $p:Q_u\twoheadrightarrow Q_t$ points in the opposite
direction and need not be simplicial on simplex boundaries. With old
blocks $\{1,2\},\{3\}$ and new singleton blocks, the object complexes
are respectively an edge disjoint from a point and a triangle boundary.
Their forward inclusion kills reduced $H_0$. But $p$ sends the new edge
$\{1,3\}$ onto all old types, a forbidden boundary face. The correct
comparison above uses a chosen section $Q_t\to Q_u$, not the reverse
quotient. Hence this is not a boundary functor on arbitrary surjections.

\begin{proposition}[Boolean expressivity does not determine Dowker shape]
Every finite simplicial complex $L$ with vertex set $Q$ can be represented
by test extents consisting of its nonempty maximal faces together with all
singletons. Its object Dowker complex is $L$, while its generated Boolean
algebra is $\Pow(Q)$. Consequently raw Dowker homotopy type cannot be
invariant under arbitrary Boolean definitional re-presentation.
\end{proposition}
\begin{proof}
All chosen extents are faces and the maximal ones generate exactly $L$
under downward closure. The singleton predicates generate $\Pow(Q)$.
Taking different complexes on the same $Q$ proves the final assertion.
\end{proof}

This does not make relational topology useless. It shows that the
\emph{designated witness family} is part of what it measures. For finite
$Q$, the Stone space of $\Pow(Q)$ is the discrete $q$-point space, not
$S^{q-2}$; these are different constructions \cite{Stone}. Erdmann's
literal-bit construction gives $S^{k-1}$ from $k$ independent bits with
all $2^k$ row types, whereas full proper Boolean saturation here gives
$S^{2^k-2}$ \cite[Section 8]{Erdmann}. The difference records the
chosen test family, not newly discovered geometry.

\section{Iterated and interacting test generation}

Let $\Gamma:\Pow(\Pow(H))\to\Pow(\Pow(H))$ be monotone and extensive.
It can represent a chosen grammar of test constructors together with
specified experimental operations. Put
\[
 \mathcal E_0=\mathcal E,\qquad
 \mathcal E_{n+1}=\Gamma(\mathcal E_n).
\]
No assumption that the generators act independently is made.

\begin{theorem}[Stable generation and first escape]\label{thm:iteration}
If $|H|=r<\infty$, the sequence stabilizes at a least
$\Gamma$-fixed family $\mathcal E_\infty$ containing $\mathcal E$ after
at most $2^r-|\mathcal E|$ strict stages. The following are equivalent:
\begin{enumerate}[label=(\roman*)]
\item $\mathcal E_\infty$ preserves the original FCA closure, residual
      equivalence, and object complex.
\item $\mathcal E_\infty\subseteq\Safe(\mathcal E)$.
\item Every test produced at every stage belongs to
      $\Safe(\mathcal E)$.
\end{enumerate}
If a structure changes at some stage, no later additive stage restores
literal equality with that original structure. The first failing stage
has a certificate of Proposition~\ref{prop:certificates} relative to the
original family.
\end{theorem}

\begin{proof}
There are $2^r$ possible extents; every strict stage adds at least one.
Monotonicity proves by induction that every fixed family containing
$\mathcal E$ contains every $\mathcal E_n$, establishing leastness.
The three equivalent conditions follow from Theorem~\ref{thm:criteria}.
Under addition, extent families and complexes can only enlarge, while
residual equivalence relations can only shrink. A previously changed
structure therefore cannot revert by further addition.
\end{proof}

\begin{remark}[Interactions matter]
It is insufficient to check only the initial generators and ignore their
compositions. On $H=\{a,b\}$, the tests $\{a\}$ and $\{b\}$ are already
present and harmless. A grammar admitting their union generates $H$.
It preserves the FCA extent closure and residual states but changes two
isolated vertices into an edge. All generated tests, not only the seeds,
must be checked. Conversely, nonempty intersections of members of
$\Safe(\mathcal E)$ stay in $\Safe(\mathcal E)$, since they remain closed
and are contained in one of the input extents.
\end{remark}

For infinite $H$, use transfinite stages with unions at limits to obtain
the least fixed point of an arbitrary monotone extensive $\Gamma$ on
the set $\Pow(\Pow(H))$. The same exact final-family criteria and
irreversibility statements apply, but there is no finite stabilization
bound. For finitary constructors, the union of the finite stages is already
closed under those constructors. Infinitary constructors need not have
this property. Neither case supplies finite termination of discovery on
an expanding, unknown state domain.

\section{Action laws and their exact domains}

\subsection{The closure lattice}
Write $L=\Ext(\mathcal E)$ and define
\[
 M_L=\{m\in M:U_m(L)\subseteq L\}.
\]
\begin{theorem}[Exact lattice action criterion]\label{thm:latticeaction}
For $m\in M$, the following are equivalent:
\begin{enumerate}[label=(\roman*)]
\item $m\in M_L$.
\item $U_m(E)\in L$ for every $E\in\mathcal E$.
\item $U_m|_L$ is a well-defined map $L\to L$ preserving arbitrary meets.
\end{enumerate}
The set $M_L$ is a submonoid. For $m\in M_L$, the map
\[
 F_m:L\longrightarrow L,\qquad F_m(A)=\cl_{\mathcal E}(\tau_m(A))
\]
is left adjoint to $U_m|_L$. On this submonoid,
\[
 U_{mn}=U_mU_n,\qquad F_{mn}=F_nF_m,\qquad U_1=F_1=\id.
\]
For arbitrary $m$, the largest subset of $L$ on which the unmodified
inverse image takes values in $L$ is
\[
 D_m^L=\{A\in L:U_m(A)\in L\}.
\]
It is closed under arbitrary intersections. The closed forward map $F_m$
is always defined, but in general only satisfies
$F_{mn}(A)\subseteq F_n(F_m(A))$.
\end{theorem}

\begin{proof}
Inverse images preserve arbitrary intersections and $H$, so checking the
column generators suffices. The formula $U_{mn}=U_mU_n$ proves the
submonoid assertion. For closed $A,B$,
\[
 F_m(A)\subseteq B\ \Longleftrightarrow\
 \tau_m(A)\subseteq B\ \Longleftrightarrow\
 A\subseteq U_m(B),
\]
proving the adjunction. Uniqueness of left adjoints proves the composition
identity. The domain assertion follows from the same intersection law.
Finally $\tau_m(A)\subseteq\cl_{\mathcal E}(\tau_m(A))$ implies the
stated forward inequality after applying $\tau_n$ and closure.
\end{proof}

\subsection{The object complex}
Put $K=K(\mathcal E)$ and
\[
 D_m^K=\{\sigma\in K:\tau_m(\sigma)\in K\},\qquad
 M_K=\{m:D_m^K=K\}.
\]
\begin{proposition}[Exact object action domain]\label{prop:objectaction}
$D_m^K$ is the maximal subcomplex on which the vertex function $\tau_m$
defines a simplicial map to $K$. The set $M_K$ is a submonoid.
For arbitrary $H$,
\[
 m\in M_K\quad\Longleftrightarrow\quad
 \forall E\in\mathcal E\ \forall\sigma\in\Fin(E)\quad
 \tau_m(\sigma)\in K.
\]
If $H$ is finite, this is equivalent to requiring that, for every old
column $E$, its image $\tau_m(E)$ is empty or is contained in an old
column. On partial domains,
\[
 \{\sigma\in D_m^K:\tau_m(\sigma)\in D_n^K\}
 \subseteq D_{mn}^K,
\]
and the two-step map equals $\tau_{mn}$ there. Equality of these domains
is not asserted: an intermediate image can leave $K$ and subsequently
return.
\end{proposition}
\begin{proof}
Faces and their images are closed under taking subsets. Every face is
contained in a column, and every finite subset of a column is a face.
This proves the domain and criterion statements. Composition of total
simplicial maps gives the submonoid assertion; direct substitution proves
the partial composition claim.
\end{proof}

\subsection{Chu identities and generated observations}
Suppose a test translation $\rho_m:T\to T$ satisfies the Chu identity
\[
 hmRt\quad\Longleftrightarrow\quad hR\rho_m(t),
 \qquad\text{equivalently}\qquad E_{\rho_m(t)}=U_m(E_t).
\tag{Chu}
\]
This implies $m\in M_L$, but it does not by itself imply $m\in M_K$.
On the test complex its exact domain is
\[
 D_m^N=\{\varnothing\}\cup
 \left\{\tau\in N(T,R):
       \left(\bigcap_{t\in\tau}E_t\right)\cap\tau_m(H)\ne\varnothing\right\}.
\]
Indeed $\bigcap_{t\in\tau}E_{\rho_m(t)}
=\tau_m^{-1}(\bigcap_{t\in\tau}E_t)$, and its nonemptiness is exactly
the displayed condition. In particular surjectivity of $\tau_m$ suffices
for a total test-side simplicial map.

\begin{proposition}[Exact partial composition for Chu actions]
Suppose Chu translations exist for $m$ and $n$. The object domains then
satisfy the stronger equality
\[
 D_{mn}^K=\{\sigma\in D_m^K:\tau_m(\sigma)\in D_n^K\}.
\]
If the test labels also satisfy $\rho_{mn}=\rho_m\circ\rho_n$, their
domains satisfy
\[
 D_{mn}^N=\{\gamma\in D_n^N:\rho_n(\gamma)\in D_m^N\}.
\]
The respective partial maps compose exactly on these domains.
\end{proposition}
\begin{proof}
Only the reverse domain inclusions require proof. If a nonempty final
object image $\tau_n\tau_m(\sigma)$ lies in a column $E_t$, then the
intermediate image $\tau_m(\sigma)$ lies in the old column
$U_n(E_t)=E_{\rho_n(t)}$. Thus it was already a face. Empty faces present
no issue. On the test side, if $h$ witnesses all tests in
$\rho_m\rho_n(\gamma)$, then $hm$ witnesses all tests in
$\rho_n(\gamma)$, by the Chu identity. The state and label composition
laws then give equality of the maps on their stated domains.
\end{proof}

\begin{example}[A Chu action need not act on either whole nerve]
Let $H=\{0,1,2\}$, and let an idempotent monoid element $a$ act by
$\tau_a(0)=\tau_a(1)=0$, $\tau_a(2)=2$. Use the columns
\[
 A=\{0,1\},\qquad B=\{2\},\qquad C=\{1,2\}.
\]
There is no universal column and every state occurs in a column.
The inverse images are $U_a(A)=A$, $U_a(B)=B$, $U_a(C)=B$, so these
define an idempotent Chu translation. Nevertheless the object edge
$\{1,2\}$ maps to the nonface $\{0,2\}$. On the test side,
$\{A,C\}$ is an edge with witness $1$, while its image $\{A,B\}$ is
not a face. The closure-lattice action is total; both nerve actions require
their stated partial domains.
\end{example}

\begin{theorem}[Action-compatible common completion]\label{thm:safeaction}
Assume $H$ is finite and nonempty. Set $S=\Safe(\mathcal E)$ and
\[
 M_S=\{m\in M:U_m(S)\subseteq S\}.
\]
Then $M_S$ is a submonoid and
\[
 m\in M_S\quad\Longleftrightarrow\quad
 U_m(E)\in S\quad\text{for every }E\in\mathcal E.
\]
The extensional completed language $S$ carries canonical test translations
$C\mapsto U_m(C)$ for $m\in M_S$, satisfying the Chu identity and the
exact composition law. Every preexisting Chu translation preserves this
completion. If $H\notin\mathcal E$, the completion still has no universal
test. The original residual quotient action, lattice action domains
$D_m^L$, and object-complex action domains $D_m^K$ are literally unchanged
by this completion.
\end{theorem}

\begin{proof}
Nonempty intersections of members of $S$ belong to $S$. If there is no
universal column, Theorem~\ref{thm:maximal} expresses each member of $S$
as a nonempty intersection of original columns. Inverse image commutes
with this intersection, so the generator criterion suffices. If there is
a universal column, the only additional case is $U_m(H)=H\in S$.
Necessity is immediate. Composition of inverse images proves the
submonoid and action laws. A preexisting Chu translation sends every
column to another old column, hence satisfies the criterion. The last
statement follows from equality of $L$, $K$, the residual relation, and
the underlying state action, which are all the data defining those domains.
\end{proof}

This action closure must be checked rather than silently imposed.
Adjoining all inverse-image tests for arbitrary operations can preserve
the residual quotient while changing the lattice or complex. Conversely,
a total lattice action or a total simplicial action is a separate
condition; neither follows from the other's definition.


\begin{proposition}[Proper closed tests on arbitrary domains]
\label{prop:proper-action}
Let $H\ne\varnothing$, $L=\Ext(\mathcal E)$, and assume $u^{-1}(L)\subseteq L$.
Then the family $L\setminus\{H\}$ is stable under $u^{-1}$ if and only if
\[
 \cl_{\mathcal E}(u(H))=H.
\]
In particular, if no old test is universal and $u$ has a Chu translation
on the old tests, all proper closed extents form a pullback-stable language,
even for infinite $H$ and infinite defining intersections.
\end{proposition}
\begin{proof}
For closed $C$, the equality $u^{-1}(C)=H$ is equivalent to $u(H)\subseteq C$.
No proper closed set contains $u(H)$ exactly when its closure is $H$.
The assumed closed-set stability supplies the other required condition.
Under a Chu translation, a proper closed extent $C$ has a nonempty
intersection representation by old columns. Its preimage is closed and is
contained in the preimage of any selected factor, itself an old proper
column. It is therefore proper.
\end{proof}

In this arbitrary-domain assertion, completion by all proper closed extents
need not be the largest simultaneous completion: on an infinite domain the
top extent may already have every finite subset witnessed without itself
being an old column. The exact largest family remains
Theorem~\ref{thm:maximal}. A label-preserving implementation may use
$T\sqcup(L\setminus\{H\})$, retaining old translations on $T$ and using
inverse image on the new extent labels. The monoid law on original labels
is separate from equality of their extents: synonymous labels may otherwise
hide a failure of the label action law.

\begin{proposition}[Preservation of existing total test actions]
Suppose $H\ne\varnothing$, no old column equals $H$, and the
instantaneous row quotient $H/{\sim_0}$ is finite. Adjoin all proper
closed extents as tests. If an old Chu translation $\rho_m$ acts
simplicially on the whole old test complex, its canonical inverse-image
translation acts simplicially on the whole completed test complex.
\end{proposition}
\begin{proof}
Every proper closed extent is an intersection of finitely many old
columns: choose, for each excluded row class, one factor in a defining
intersection that excludes it. There are finitely many classes, and
properness makes the chosen family nonempty. Any nonempty intersection
of finitely many completed test extents is consequently a nonempty
finite intersection of old columns. Totality of the old test-side map
says that the inverse image of such an intersection is nonempty. This
is exactly the face criterion for the completed test-side map.
The completed inverse-image tests themselves are proper closed extents
by the Chu identity and closure under nonempty intersections.
\end{proof}

The finite-row-quotient assumption concerns this preservation of total
test-side actions, not the unconditional preservation of the object
complex or its partial domains. For an infinite defining conjunction,
nonemptiness of every finite inverse-image intersection need not imply
nonemptiness of the full inverse-image intersection. No equivariance of
the homotopy retraction in the next section follows from this proposition.


\section{Literal complexes and canonical homotopy preservation}

In this section abbreviate $K=K(\mathcal E)$ and $N=N(T,R)$;
superscript $+$ denotes the corresponding complexes after adjoining the
specified new labeled tests. All abstract simplicial realizations use the
weak CW topology.
A newly labeled nonempty test always adds a test-side vertex, even when its
extent is a duplicate. Thus exact equality of the labeled test complexes
holds precisely when every added extent is empty. This is stronger than,
and different from, exact equality of the object complexes.


\begin{theorem}[Labeled test-side homotopy under domination]\label{thm:testnerve}
Adjoin labeled tests $g\in G$ of extents $C_g$. Suppose every nonempty
$C_g$ is contained in some old extent $E_{d(g)}$. Then the inclusion
\[
 i:N(T,R)\longrightarrow N(T\sqcup G,R^+)
\]
is a homotopy equivalence on geometric realizations. An explicit simplicial
retraction sends each old vertex to itself and each new vertex $g$ to
$d(g)$. Thus every object-complex-preserving extension on finite $H$
has this property. The test-side complex can have new vertices and faces;
it is not claimed to be literally equal to the old complex.
\end{theorem}
\begin{proof}
An added empty-extent label is not a vertex and needs no vertex image.
Choose $d(g)$ for the other labels. A witness for any finite collection of
old and new tests also witnesses its image under the proposed retraction,
because each replacement extent contains the original one. This proves
simpliciality. The retraction $r$ satisfies $ri=\id$. For every simplex
$\tau$, the union $\tau\cup ir(\tau)$ has the same witness, so $ir$ is
contiguous to the identity. Contiguous simplicial maps have homotopic
realizations.
\end{proof}

The choices $d(g)$ need not commute with a Chu translation. The theorem
does \emph{not} claim an equivariant retraction or an equivariant homotopy.
For infinite $H$, equality of the object complex supplies only the
finite-subset criterion of Theorem~\ref{thm:criteria}; a dominating old
column need not exist. The functorial Dowker theorem still compares the
two test complexes up to homotopy through the unchanged object complex
\cite{BrunSalbu}, but this particular vertex-choice retraction is not then
available in general.


\begin{corollary}[Exact history preservation always preserves nerve homotopy]
\label{gt:general-homotopy}
For arbitrary sets, if the object-complex criterion in Theorem~\ref{thm:criteria} holds, the
canonical inclusion $N\hookrightarrow N^+$ is a homotopy equivalence,
even if no global dominators $t_j$ exist. More generally the canonical
inclusion $K\hookrightarrow K^+$ is a homotopy equivalence if and only if
the canonical inclusion $N\hookrightarrow N^+$ is one.
\end{corollary}
\begin{proof}
The context inclusion $(\mathrm{id}_H,T\hookrightarrow T\sqcup J)$ is
an incidence-preserving morphism. Functorial Dowker duality gives a
commutative square in the homotopy category with these two inclusions
as its horizontal arrows and homotopy equivalences as its vertical
arrows. Therefore either horizontal arrow is invertible if and only
if the other is. This version for arbitrary sets and finite simplices
is supplied, for example, by the strictly natural rectangle-complex
span in Brun--Salbu~\cite[Theorems 5.1--5.2]{BrunSalbu}.
\end{proof}


\subsection{Dominated completion on exact sequential action domains}

The domination proof can preserve more than an isolated homotopy type.
Suppose the incidence presentation carries a history operation
$u:H\to H$ and a test operation $v:T\to T$ satisfying
\[
 E_{v(t)}=u^{-1}(E_t).
\]
Assume the extended vocabulary $T^+=T\sqcup J$ carries an operation
$v^+:T^+\to T^+$ extending $v$ and satisfying the same inverse-image
law for every old or new extent. Write $E_e$ for the extent of any label
$e\in T^+$, and $D_j=E_j$ for a new label $j\in J$. For a subset
$S\subseteq H$, let
$N_S$ and $N_S^+$ be the two test nerves formed using witnesses only
in $S$. The exact domain of the test operation is
$N_{u(H)}$, respectively $N^+_{u(H)}$: a face is in this domain
exactly when all its pulled-back tests have a common witness.

\begin{theorem}[Domainwise comparison and homotopy commutation]
\label{gt:action-domains}
Under the domination hypotheses of Theorem~\ref{thm:testnerve}, the
chosen retraction $r:N^+\to N$ restricts, for every $S\subseteq H$,
to a retraction $r_S:N_S^+\to N_S$ contiguous to the identity after
inclusion. In particular it is an equivalence on the exact $u$-domains.
The two simplicial maps
\[
 r\,v^+,\quad v\,r_{u(H)}:N^+_{u(H)}\longrightarrow N
\]
are contiguous, relative to the old domain $N_{u(H)}$.
The forward inclusions commute with the test actions strictly. Thus
these inclusions and their domainwise homotopy inverses commute with
the actions in the homotopy category.
\end{theorem}
\begin{proof}
Intersecting $D_j\subseteq E_{t_j}$ with $S$ preserves the containment.
A new vertex actually present in $N_S^+$ has a witness in $D_j\cap S$,
so its chosen old dominator is actually present in $N_S$. The proof of
Theorem~\ref{thm:testnerve} works with that same witness in $S$, including
when the restricted nerves have no vertices.

Now let $\tau$ be a nonempty face of $N^+_{u(H)}$. Choose $x\in H$
with $u(x)$ lying in every extent indexed by $\tau$. For every
$e\in\tau$, inverse-image compatibility gives $x\in E_{v^+(e)}$.
The chosen retraction either fixes this test or replaces it by a
dominating old extent, so $x\in E_{r(v^+(e))}$. On the other hand,
$E_e\subseteq E_{r(e)}$ gives $u(x)\in E_{r(e)}$, and therefore
$x\in E_{v(r(e))}$. Thus $x$ witnesses the union of the two image
simplices, which proves contiguity. All maps are well-defined on
actual vertices: the same argument shows that their image extents
are nonempty. On old tests both formulas equal $v$, so the homotopy
is relative to the old domain. Strict commutation of forward
inclusions follows from the assumption that $v^+$ extends $v$.
\end{proof}

Apply this theorem separately to every word $w$ in an acted context,
using $u_w$ and its contravariant test operation $v_w$. The domination
choices can remain fixed, and the argument supplies equivalences on
every exact word-domain and commuting squares in the homotopy
category. No unproved strictly equivariant retraction or system of
higher coherent homotopies is asserted. If history-side partial
operations use the domains $\{\sigma\in K:u_w(\sigma)\in K\}$,
then equality $K^+=K$ preserves those domains and the maps on them
literally, since the history operations themselves are unchanged.


\subsection{An exact criterion for one additional test}

The next theorem strengthens the common sufficient test that the
attachment intersection be contractible. It is a direct CW attachment
calculation using ordinary integral homology, van Kampen, covering spaces,
and Whitehead's theorem. No novelty is claimed. It concerns the
\emph{canonical inclusion}, not merely an abstract coincidence of the
two spaces' homotopy types.

\begin{theorem}[Exact one-simplex attachment criterion]
\label{gt:one-test}
Let $K$ be a finite simplicial complex, with possibly some unused ambient
vertices, and let $D$ be a finite nonempty subset of that ambient vertex
set. Put $B=\Delta(D)$, $K^+=K\cup B$, and $A=K\cap B=K[D]$.
The inclusion $|K|\hookrightarrow|K^+|$ is a homotopy equivalence if and
only if both of the following hold:
\begin{enumerate}
\item $|A|$ is nonempty and integrally acyclic:
$\widetilde H_n(|A|;\mathbb Z)=0$ for every $n\geq0$.
\item For a basepoint $a\in|A|$, the inclusion-induced homomorphism
\[
 \pi_1(|A|,a)\longrightarrow\pi_1(|K|,a)
\]
has trivial image.
\end{enumerate}
Item 1 implies that $|A|$ is path connected, so the choice of basepoint
does not affect triviality of the image. Only the component of $|K|$
containing $A$ is involved; all other components are unchanged.

The identical assertion holds without finiteness assumptions for
abstract simplicial complexes with finite simplices and weak CW
realizations. It also holds for a union $X\cup_A B$ of CW subcomplexes
when $B$ is any nonempty contractible CW complex.
\end{theorem}
\begin{proof}
Spaces are written in place of their realizations inside the proof. Simplicial
subcomplexes are CW subcomplexes, so their inclusions are cofibrations.

\emph{Necessity.} If $A$ has no points, the newly attached nonempty
simplex is a disjoint additional component. The inclusion cannot be a
homotopy equivalence. Otherwise the relative simplicial chain complexes
are naturally isomorphic:
\[
 C_*(K^+,K;\mathbb Z)\cong C_*(B,A;\mathbb Z).
\]
Indeed, on both sides the basis in each degree consists of exactly the
faces of $B$ not already in $A$, with the same relative boundary maps.
Since $B$ is contractible and $A\neq\varnothing$, the relative long
exact sequence gives, for every $n\geq1$,
\begin{equation}\label{gt:relative-homology}
 H_n(K^+,K;\mathbb Z)\cong H_n(B,A;\mathbb Z)
                       \cong\widetilde H_{n-1}(A;\mathbb Z).
\end{equation}
A homotopy equivalence has zero relative homology, so $A$ is acyclic.
In particular its reduced zeroth homology vanishes, making the nonempty
CW complex $A$ path connected. The composite
$A\hookrightarrow K\hookrightarrow K^+$ factors through the
contractible complex $B$, so it induces the zero homomorphism on
$\pi_1$. Since $\pi_1(K,a)\to\pi_1(K^+,a)$ is an isomorphism,
the homomorphism $\pi_1(A,a)\to\pi_1(K,a)$ must already have been zero.

\emph{Sufficiency: components and fundamental groups.} Let $X$ be the
component of $K$ containing the connected nonempty $A$, and let
$Y=X\cup_A B$. The other components of $K$ are left unchanged, so it
suffices to prove $X\hookrightarrow Y$ a homotopy equivalence.
The van Kampen theorem for a union of connected CW subcomplexes gives
\[
 \pi_1(Y,a)\cong\pi_1(X,a)
       /\langle\!\langle\operatorname{im}\pi_1(A,a)\rangle\!\rangle.
\]
Here $B$ has trivial fundamental group. By item 2 the normal subgroup
is trivial, hence $\pi_1(X,a)\to\pi_1(Y,a)$ is an isomorphism.

\emph{Universal-cover calculation.} Put $G=\pi_1(Y,a)$ and take a
universal cover $p:\widetilde Y\to Y$. The inverse image
$\widetilde X=p^{-1}(X)$ is connected because $\pi_1(X)\to G$ is
surjective, and simply connected because that homomorphism is injective.
Thus $\widetilde X$ is the universal cover of $X$.
After fixing a lift of $a$, the inverse image of the simply connected
$B$ is a disjoint union of copies $B_g$ indexed by $G$. Each copy
meets $\widetilde X$ in a copy
$A_g$ of $A$: the restriction of $p$ to $B_g$ is a homeomorphism onto
$B$, and its inverse image of $A$ is exactly this intersection. All
cells of $\widetilde Y\setminus\widetilde X$ lie in these copies.
Therefore
\begin{equation}\label{gt:cover-relative}
 C_*^{\mathrm{cell}}(\widetilde Y,\widetilde X;\mathbb Z)
       \cong\bigoplus_{g\in G}C_*^{\mathrm{cell}}(B,A;\mathbb Z).
\end{equation}
Every summand is acyclic by item 1 and the relative long exact sequence.
Direct sums of chain complexes preserve this acyclicity. Hence
$\widetilde X\hookrightarrow\widetilde Y$ induces isomorphisms on
all integral homology groups. Both spaces are simply connected CW
complexes; the homological form of Whitehead's theorem makes this lift
a homotopy equivalence. Covering maps identify higher homotopy groups,
so $X\hookrightarrow Y$ induces isomorphisms on $\pi_n$ for $n\geq2$,
as well as on $\pi_1$ and components. Whitehead's theorem now proves
the claimed homotopy equivalence. The standard covering and Whitehead
facts used here are in Hatcher~\cite[Sections 1.3 and 4.1,
Corollary 4.33]{Hatcher}.

No step uses finiteness except to avoid infinite chain notation. For
arbitrary CW subcomplexes the relative cellular chain argument and the
direct sum in \eqref{gt:cover-relative} are the same. In the stated
simplicial setting, all realizations and covering spaces are CW spaces.
\end{proof}

\begin{corollary}[Useful sufficient and insufficient checks]
\label{gt:checks}
For one nonempty added test with extent $D$:
\begin{enumerate}
\item If $|K[D]|$ is nonempty and contractible, its addition preserves
the canonical homotopy type. In particular this holds if $K[D]$ has a
cone vertex or is certified collapsible.
\item Integral acyclicity of $K[D]$ is necessary but not sufficient
in general. It is sufficient when $K[D]$ lies in a simply connected
component of $K$.
\item If $K[D]$ is disconnected, or has any nonzero reduced integral
homology, the canonical inclusion cannot be a homotopy equivalence.
Nonzero rational homology is an easy obstruction, but zero rational
homology alone is insufficient because it can miss torsion.
\end{enumerate}
\end{corollary}

There is also a direct test-side expression. For a single new test $j$
of extent $D$, its old-vertex link is
\[
 L_D=\{\varnothing\}\cup
 \{\varnothing\neq\tau\subseteq T:\tau\text{ finite and }
                         D\cap\bigcap_{t\in\tau}E_t\neq\varnothing\}.
\]
Thus $N^+=N\cup(j*L_D)$ and $N\cap(j*L_D)=L_D$.
Theorem~\ref{gt:one-test}, in its contractible-CW version, applies with
$L_D$ in place of $A$. The restricted incidence relation on $D\times T$
has history complex $K[D]$ and test complex $L_D$, so Dowker duality
identifies their homotopy types. Naturality relates the attachment maps
to the old complexes, not merely these two abstract spaces. This explains
why the exact criterion agrees on both sides.

\begin{example}[Why the fundamental-group clause is needed]
There is a finite connected acyclic simplicial complex $A_0$ with
$\pi_1(A_0)\neq1$. One explicit finite CW model starts with two loops
$a,b$ and attaches two $2$-cells with words $a^5b^{-3}$ and
$b^3(ab)^{-2}$. The cellular boundary matrix is
\[
 \begin{pmatrix}5&-2\\-3&1\end{pmatrix},\qquad\det=-1,
\]
so reduced integral homology vanishes. The fundamental group is
nontrivial: send $a$ to $(1\ 2\ 3\ 4\ 5)$ and $b$ to $(1\ 4\ 2)$
in the alternating group on five letters. Both relations hold because
$a^5=b^3=(ab)^2=1$, while the image of $a$ is nonidentity. The standard
finite simplicial replacement of a finite CW complex gives $A_0$ with
these same homotopy invariants. This model is the classical acyclic
example discussed by Hatcher~\cite[Example 2.38 and Theorem 2C.5]{Hatcher}.

Take old test extents to be the maximal faces of $A_0$, so that the
old history complex is exactly $K=A_0$. Add truth on its finite vertex
set. The new history complex is a full simplex. Homology is preserved,
but homotopy is not: the nontrivial fundamental group is killed. In
Theorem~\ref{gt:one-test}, the map on $\pi_1(A_0)$ is the identity,
not the trivial homomorphism.

Conversely let $K=v*A_0$ be a cone, put $D=V(A_0)$, and again use the
maximal faces of $K$ as old extents. Then $K[D]=A_0$ is not contractible,
but $K$ is simply connected and $A_0$ is acyclic. The theorem proves
that adjoining $D$ is a homotopy equivalence. Hence contractibility of
the attachment intersection is sufficient, but genuinely not necessary.
\end{example}

\subsection{Several tests: safe stages, exact global tests, and cautions}

For a finite ordered list $D_1,\ldots,D_m$, put
$K_0=K$ and $K_i=K_{i-1}\cup\Delta(D_i)$. At stage $i$, use the
\emph{current} intersection $A_i=K_{i-1}[D_i]$, not just the
intersection with $K_0$. If every nonempty attachment passes
Theorem~\ref{gt:one-test}, then $K_0\hookrightarrow K_m$ and
$N_0\hookrightarrow N_m$ are homotopy equivalences. For arbitrary
families dominated by old extents, Theorem~\ref{thm:testnerve} already
provides a single simultaneous certificate and needs no ordering.

Stagewise success is sufficient but not necessary for success of the
final inclusion in a chosen order. For instance, start with the path
$\{1,2\}\cup\{2,3\}$, first adjoin $D_1=\{1,3\}$, then
$D_2=\{1,2,3\}$. The first stage creates a circle and the second
fills it, so neither stage inclusion is a homotopy equivalence; the
initial-to-final inclusion is nevertheless an equivalence between
contractible spaces. This also shows why one must test the map and
the intended stopping stage rather than only count the added columns.

For the general finite final inclusion $j:|K|\hookrightarrow|K^+|$,
an exact global criterion is that $j$ be bijective on components, induce
isomorphisms on fundamental groups of corresponding components, and
that the induced maps between their universal covers be integral
homology isomorphisms. Necessity follows from lifting a homotopy
equivalence, and sufficiency is the universal-cover Whitehead argument
above. Unlike the one-simplex case, these cover-relative chain complexes
need not split into independent attachments. This is a characterization,
not an efficient general algorithm: neither arbitrary fundamental-group
questions nor arbitrary homotopy equivalence should be advertised as
routine finite computation.

For finite incidence tables, exact preservation \emph{is} elementary:
compare each nonempty new extent with the old extents and look for a
superset. With $n=|H|$, $m=|T|$, and $q$ added tests, a direct bit-table
implementation needs at most $O(nmq)$ Boolean comparisons. A successful
comparison supplies the explicit retraction witness as well as the
equality certificate; enumeration of all simplices is unnecessary.

\subsection{Observation covers and the meaning of a topology change}

Suppose now that $H$ is a sample in a space $Z$, with tests representing
membership in subsets $U_t\subseteq Z$, so $E_t=H\cap U_t$.
The sample nerve has face $\tau$ precisely when
$H\cap\bigcap_{t\in\tau}U_t\neq\varnothing$. It equals the true
cover nerve exactly when every nonempty finite cover intersection has
a sample witness. Formula generation alone does not prove that sampling
condition, and finite-sample identities do not prove identities on all
of $Z$.

\begin{proposition}[A cover-preserving interpretation]
\label{gt:cover}
Let $Z$ be paracompact Hausdorff and let $\mathcal U$ be a finite good
open cover: every nonempty finite intersection of members is
contractible. Let $\mathcal V$ be a finite family of additional open
subsets of $Z$. If $\mathcal U\cup\mathcal V$ is also a good cover,
then the canonical nerve inclusion
\[
 |N(\mathcal U)|\hookrightarrow|N(\mathcal U\cup\mathcal V)|
\]
is a homotopy equivalence. In particular one may adjoin any family of
nonempty finite intersections of original cover elements. If a sample
is nerve-complete for the enlarged cover, the same conclusion applies
to its observed test nerve.
\end{proposition}
\begin{proof}
The cover inclusion is a morphism of covers over the identity map of
$Z$. The functorial good-cover nerve theorem identifies the induced
nerve map, up to homotopy equivalences on both sides, with
$\mathrm{id}_Z$. Therefore it is a homotopy equivalence; a precise
natural-span treatment is in
Bauer--Kerber--Roll--Rolle~\cite{BauerEtAl}. Adding finite intersections
preserves openness, does not change the union, and makes every new
finite intersection itself an intersection of original members, hence
empty or contractible. Nerve-completeness identifies the observed
complexes with the true nerves. In the finite-intersection case the
combinatorial domination theorem independently proves the inclusion
equivalence without any good-cover assumptions.
\end{proof}

Arbitrary unions behave differently. Even when the original cover is
good, a union may be noncontractible or may have disconnected or
noncontractible intersections with other members. Adding $Z$ itself
always cones the nerve; if $Z$ is noncontractible, the enlarged cover
is not good because one of its members is $Z$. A logical truth test
therefore cannot be treated as another harmless member of a good
geometric cover. For example, start with a finite good open arc cover of $S^1$. Its
complete union is $S^1$; treating that finite disjunction as one extra
cover member cones the nerve. The new cover fails the good-cover
hypothesis because the added member is the noncontractible circle.
This is not a counterexample to the nerve theorem. Similarly, infinite
intersections of open sets need not be open. Their combinatorial domination remains valid, but the
good-open-cover theorem no longer applies automatically.

There are two legitimate outcomes. A generated predicate may preserve
the selected topological observable, certified above, or it may alter
that observable because a different observation family is now being
studied. The latter is not intrinsically a harmful representation move.
It becomes an error when the changed nerve is asserted to encode the
same invariant without the needed preservation theorem, or when a
topology inferred from one cover is misidentified with the topology of
an unchanged ambient space. Native test generation can be useful even
when it intentionally changes topology; its effect must be stated at
the chosen level of structure.


\section{A scientific instance and its assumptions}
\label{sec:oscillator}

Consider a harmonic oscillator with known $\omega>0$,
\[
 \dot x=v,\qquad \dot v=-\omega^2x,
 \qquad H=\mathbb R^2,
\]
and instantaneous real-valued probe $p(x,v)=x$. States with the same
position and different velocity agree on every pointwise function of $p$,
including every Boolean combination of position-threshold tests. Their
future positions need not agree. The known flow gives the pulled-back probe
\[
 p_\tau(x,v)=p(\Phi_\tau(x,v))
     =x\cos(\omega\tau)+\frac{v}{\omega}\sin(\omega\tau).
\]
When $\sin(\omega\tau)\ne0$, the two probes $p,p_\tau$ identify the
whole state:
\[
 x=p,\qquad
 v=\omega\,\frac{p_\tau-p\cos(\omega\tau)}{\sin(\omega\tau)}.
\]
At the quarter-period $\tau=\pi/(2\omega)$, $p_\tau=v/\omega$.
Equivalently, the Lie derivatives of $p$ satisfy
$L_fp=v$, $L_f^2p=-\omega^2p$; their linear span closes after the
position and velocity probes. This is a concrete instance of observability,
not just renaming a binary column. It belongs to the established program
of generating observation functions from dynamics \cite{HermannKrener}.

There is an exact sampling obstruction. If the only allowed time step is
$\tau=k\pi/\omega$, then
\[
 p(\Phi_{n\tau}(x,v))=(-1)^{kn}x
 \qquad(n\ge0).
\]
All such observations remain independent of velocity. The full residual
quotient for that sampling protocol identifies every velocity at a fixed
position, whereas a nonresonant time step separates them. The action system,
not merely the algebra of formulas, determines which distinctions can emerge.

The exact recovery is also ill-conditioned near a resonant step. If the
position readings have bounded errors $|e_0|\le\eta_0$ and
$|e_\tau|\le\eta_\tau$, the displayed reconstruction obeys
\[
 |\widehat v-v|
 \le\frac{\omega\bigl(\eta_\tau+
               |\cos(\omega\tau)|\eta_0\bigr)}
              {|\sin(\omega\tau)|}.
\]
This follows by substituting the perturbed readings and applying the
triangle inequality. It illustrates why semantic distinguishability alone
is not a robustness or sample-complexity guarantee.

The example assumes the dynamical law, frequency, timing, and ability to
measure position along the same trajectory. It does not infer that law
from the original position-only incidence table. If $\omega$ is unknown,
two position readings generally do not identify both frequency and velocity.
Finite thresholds do not give exact real-valued reconstruction, and replacing
the state space by a smooth quotient requires additional regularity. Thus
new operational characterizations can expose a physically useful state
variable, while their evidential and modeling assumptions remain explicit.

\section{Incomplete evidence and possible completions}
\label{sec:evidence}

All binary incidence theorems above concern an exact relation on a declared
domain. A blank entry may be interpreted as false only when falsity is
actually the intended semantics. Untested, unrecorded, or missing observations
must not silently become negative results. Even a complete finite table
need not describe every possible object or justify a population law.

For each observed $h$, let $\mathcal V_h\ne\varnothing$ be the set of
valuations compatible with the record and the explicitly adopted background
constraints. For a formula $\varphi$, define
\[
 \operatorname{must}(\varphi)
 =\{h:\forall\nu\in\mathcal V_h,\ \nu\models\varphi\},
 \qquad
 \operatorname{may}(\varphi)
 =\{h:\exists\nu\in\mathcal V_h,\ \nu\models\varphi\}.
\]
Known truth is membership in $\operatorname{must}$, known falsity is
nonmembership in $\operatorname{may}$, and the remaining case is unknown.
These bounds are exact relative to the chosen compatible valuations, not
an assurance that the chosen possibilities contain the actual world.

\begin{proposition}[Completion semantics retains correlations]
The following identities and inclusion hold:
\begin{align*}
 \operatorname{must}(\neg\varphi)&=H\setminus\operatorname{may}(\varphi),&
 \operatorname{may}(\neg\varphi)&=H\setminus\operatorname{must}(\varphi),\\
 \operatorname{must}(\varphi\wedge\psi)
 &=\operatorname{must}(\varphi)\cap\operatorname{must}(\psi),&
 \operatorname{may}(\varphi\wedge\psi)
 &\subseteq\operatorname{may}(\varphi)\cap\operatorname{may}(\psi).
\end{align*}
The last inclusion can be strict. Independent lower and upper sets therefore
need not suffice for exact evaluation of generated tests.
\end{proposition}
\begin{proof}
The first three equalities follow by distributing the universal quantifier
and exchanging existential and universal quantifiers under negation.
A common valuation witnessing a conjunction witnesses both conjuncts,
proving the inclusion. With one unknown atom $p$, both $p$ and $\neg p$
may hold in different completions, but $p\wedge\neg p$ holds in none.
\end{proof}

Strong Kleene three-valued evaluation is another legitimate semantics, but
it is not the same as this supervaluational completion semantics: with
unknown $p$, it can leave $p\vee\neg p$ unknown, whereas every classical
completion makes that formula true. The evidence convention must therefore
be stated before applying Boolean test-generation rules.

For a concrete error, take two observed objects with
$p(h_1)=1$, $q(h_1)$ unknown, $p(h_2)=0$, $q(h_2)=1$.
Filling the blank with zero makes $p\Rightarrow\neg q$ appear valid.
The compatible completion $q(h_1)=1$ refutes it. Attribute exploration with
partial descriptions has a substantial existing theory
\cite{BorchmannGanter}; the exact-table lemmas here do not replace it.

\begin{proposition}[An evidence obstruction]\label{prop:evidence}
Let $W$ be a set of possible worlds and $e:W\to D$ the available record.
If $e(w)=e(w')$ but a proposition $\varphi$ has different truth values
in $w$ and $w'$, no function of that record alone can assign a definite
truth value to $\varphi$ soundly throughout $W$.
\end{proposition}
\begin{proof}
The function receives identical input in both worlds and must return the
same value, which is wrong in one of them.
\end{proof}

New logical consequences can be warranted because they hold in all
compatible worlds. New measurements can shrink that set; new assumptions
can also shrink it, but give conditional conclusions. These are different
ways of learning or reorganizing knowledge. A generated test should pass
three separate checks: its meaning is defined; an admissible procedure can
execute it at stated precision and cost; and the evidence warrants its
answer or statistical law. Mathematical measurability does not imply
computability, and an infinite ``eventually'' test can encode a halting
problem even when each finite transition is executable.

\subsection{Prediction does not automatically give a Markov quotient}

The deterministic congruence theorem must not be transferred unchanged to
stochastic systems by replacing truth with equality of future output laws.
Consider states $x,y,u,v,w,a,b$, with output $1$ only at $a$. Let $a,b$
be absorbing and let
\[
 u\to a,\quad v\to b,\quad
 w\to\tfrac12a+\tfrac12b,\quad
 x\to\tfrac12u+\tfrac12v,\quad y\to w.
\]
The output processes from $x,y$ have the same law: two initial zeros,
then an equal mixture of an all-one and an all-zero tail. But $u,v,w$
have distinct future-output laws, and the successor distributions from
$x,y$ over those equivalence classes differ. The trace quotient is not
Markov lumpable. For an exact Markov quotient one needs closure under
probability modalities such as
$U\mapsto\{h:K_a(h,U)>r\}$, not just comparison of output traces.
For predictive states of observed histories, conditional-law and almost-sure
qualifications are different again \cite{ShaliziCrutchfield}.

\section{Theory extension beyond a fixed incidence domain}
\label{sec:limits}

A fixed-domain completion is useful but does not capture all mathematical
or scientific theory formation. A theory can change its signature, introduce
new sorts, construct quotients and limits, compare models, and organize
proofs. These require syntax, model categories, and satisfaction compatible
with translation, as in categorical logic and institutions
\cite{Lawvere,GoguenBurstall}. The membership table of models and sentences
forgets model morphisms, substitution, contexts, and proof structure.

A definitional extension adds a new predicate $P$ with
$P(\mathbf x)\leftrightarrow\varphi(\mathbf x)$ in the old language.
Every old model has a unique expansion, and substitution eliminates $P$,
so old-language consequences are unchanged. Function definitions need
existence and uniqueness; recursive or overloaded definitions need their
own hypotheses \cite{Xue,GengelbachWeber}. This does not imply preservation
of a chosen FCA lattice or nerve formed by treating every newly named
formula as a primitive test. The earlier counterexamples show why.
An operational extension adds a new admissible experiment or construction;
a substantive theory extension can instead restrict the admitted models
by adding hypotheses not entailed before. Neither is justified merely by
closure of the old syntax.

Even a limit of increasingly discriminating descriptions need not consist
only of actual objects. Let $H$ be the eventually-zero binary sequences,
and let stage $n$ observe the first $n$ bits. Each finite quotient contains
all $n$-bit strings, but their inverse limit is the space of \emph{all}
infinite binary sequences. The canonical map from the fully distinguished
$H$ into this inverse limit is injective and not surjective. Ideal completion
has added possible profiles, not discovered members of the original domain.
Compactness or other realizability assumptions can change this conclusion,
but must be supplied.

Classifying toposes give a different, richer universal construction. For a
geometric theory $\mathbb T$, its classifying topos $\mathcal E_{\mathbb T}$
represents models across Grothendieck toposes by
\[
 \operatorname{Geom}(\mathcal F,\mathcal E_{\mathbb T})
 \simeq \mathbb T\text{-}\operatorname{Mod}(\mathcal F).
\]
This is not the observational quotient of one dataset. Transfer between
presentations requires a proved interpretation or equivalence and a chosen
invariant; it does not automatically preserve literal syntax or proof
complexity. It is a useful broader direction, not a consequence of the
incidence calculations here \cite{Caramello}.

For a substantive theory of generated characterizations, the next theorem
would therefore have to specify admissible changes of signature and
experimental access, the intended consumers and their preservation levels,
and identifiability or evidence conditions. The present results supply a
precise fixed-domain baseline and several obstructions. They do not rank
scientific hypotheses, invent physical procedures, validate an unknown
law, or guarantee discovery beyond the supplied language and operations.

\section{Established results and the contribution boundary}

The extent criteria are standard FCA reduction expressed as an extension
problem: a conjunction-defined attribute is reducible, and adding it preserves
the concept hierarchy \cite{GanterWille,GanterRudolphStumme}. Infomorphisms,
closure spaces, and induced lattice adjunctions already have substantial
functorial treatments \cite{ShenZhang,KroetzschEtAl}. Those morphisms must
not be confused with covariant incidence-preserving pairs, the morphisms
used in the functorial Dowker comparison.

The greatest invariant quotient and future-depth refinement belong to
Moore--Myhill--Nerode observability and universal coalgebra
\cite{Moore,Rutten}. Ranzato--Tapparo's complete abstractions and strong
preservation provide a particularly close and richer baseline for refining
representations to respect an inductively generated observation language,
including nonpartitioning abstractions \cite{RanzatoTapparo}. Angluin's
learning theorem and coalgebra learning via duality address actual learning
under their query, finiteness, and structural assumptions
\cite{Angluin,BarloccoKupkeRot}. The latter explicitly uses formulas as
tests and proves correctness and termination. Expressivity of a chosen
coalgebraic logic requires hypotheses such as separation by predicate
liftings; it is not implied by merely adjoining Boolean syntax
\cite{Pattinson}.

Dowker duality, natural comparison spans, and nerve theorems supply the
topological foundations \cite{ChowdhuryMemoli,BrunSalbu,BauerEtAl}.
The one-simplex criterion proved here is a specialization of standard
CW attachment, covering-space, and Whitehead theory \cite{Hatcher}.
Erdmann already studies the topology of logical attribute transformations,
including literal-bit spheres, Boolean OR changes, the risk of generating
truth, and the need for physically observable attributes
\cite[Sections 8--9]{Erdmann}. The full disjunction and punctured-saturation
statements make particular consequences exact; the underlying warning is
not new.

Green et al. explicitly integrate hypergraph intersection closures, FCA,
Dowker complexes, and cosheaves \cite{GreenEtAl}. Their Definition 25
retains $(G,M,L,\gamma,\mu)$, distinguishing an enriched concept lattice
from a bare complete lattice. Their Definition 24 also uses a strong context-morphism condition:
$f(A)'=g(A')$ and $g(B)'=f(B')$ for all subsets, rather than arbitrary
covariant incidence-preserving maps or Chu infomorphisms. Their reconstruction
statements must not be
read as invariance under bare lattice isomorphism or under freely adding
truth; object and attribute labels and their maps retain additional incidence
information. Their diagram also need not commute. Thus neither representing
one relation in several ways nor adding a category-theoretic vocabulary is
by itself a new unification theorem.

The defensible output is the exact joint preservation problem, its maximal
completion, its dynamic and epistemic qualifications, and a collection of
proved separation results. No claim of priority is made for their combined
formulation. A new scientific or phenomenological theory would need a further
result linking licensed operations and evidence to a specified target, rather
than identifying every formula-generated topology with the structure of that
target.

\begin{thebibliography}{99}
\bibitem{GanterWille}
B. Ganter and R. Wille, \emph{Formal Concept Analysis: Mathematical
Foundations}, Springer, 1999.
\url{https://doi.org/10.1007/978-3-642-59830-2}.

\bibitem{GanterRudolphStumme}
B. Ganter, S. Rudolph, and G. Stumme,
\emph{Explaining Data with Formal Concept Analysis}, 2019,
Section 3.11, Definition 15 and Proposition 8.
\url{https://iccl.inf.tu-dresden.de/w/images/4/49/IntroFCA-RW2019.pdf}.

\bibitem{ShenZhang}
L. Shen and D. Zhang, \emph{The concept lattice functors},
International Journal of Approximate Reasoning 54(1) (2013), 166--183.
\url{https://doi.org/10.1016/j.ijar.2012.07.002}.

\bibitem{KroetzschEtAl}
M. Kr\"otzsch, P. Hitzler, and G.-Q. Zhang,
\emph{Morphisms in Context}, 2005.
\url{https://iccl.inf.tu-dresden.de/w/images/1/1f/ICCS05_context-morphisms_long.pdf}.

\bibitem{BorchmannGanter}
D. Borchmann and B. Ganter,
\emph{Abstract Attribute Exploration with Partial Object Descriptions}, 2015.
\url{https://arxiv.org/abs/1511.06191}.

\bibitem{Moore}
E. F. Moore, \emph{Gedanken-Experiments on Sequential Machines},
in C. E. Shannon and J. McCarthy (eds.), \emph{Automata Studies}, 1956.
\url{https://doi.org/10.1515/9781400882618-006}.
\url{https://www.cs.cmu.edu/~cdm/resources/Moore1956-gedanken-experiments.pdf}.

\bibitem{Rutten}
J. J. M. M. Rutten, \emph{Universal coalgebra: A theory of systems},
Theoretical Computer Science 249 (2000), 3--80.
\url{https://ir.cwi.nl/pub/48/0048D.pdf}.

\bibitem{RanzatoTapparo}
F. Ranzato and F. Tapparo,
\emph{Generalizing the Paige--Tarjan Algorithm by Abstract Interpretation}.
\url{https://www.math.unipd.it/~ranzato/papers/ArXiV0612120.pdf}.
See also the authors' strong-preservation treatment:
\url{https://www.math.unipd.it/~ranzato/papers/esop04.pdf}.

\bibitem{Angluin}
D. Angluin, \emph{Learning regular sets from queries and counterexamples},
Information and Computation 75 (1987), 87--106.
\url{https://doi.org/10.1016/0890-5401(87)90052-6}.
\url{https://homepages.math.uic.edu/~lreyzin/papers/angluin87.pdf}.

\bibitem{BarloccoKupkeRot}
S. Barlocco, C. Kupke, and J. Rot,
\emph{Coalgebra Learning via Duality}, FoSSaCS 2019, 62--79,
Theorems 33 and 37.
\url{https://doi.org/10.1007/978-3-030-17127-8_4}.
\url{https://arxiv.org/abs/1902.05762}.

\bibitem{Pattinson}
D. Pattinson, \emph{Expressive Logics for Coalgebras via Terminal Sequence
Induction}, Notre Dame Journal of Formal Logic 45(1) (2004), 19--33.
\url{https://users.cecs.anu.edu.au/~dpattinson/Publications/ndjfl2004.pdf}.

\bibitem{ChowdhuryMemoli}
S. Chowdhury and F. M\'emoli,
\emph{A functorial Dowker theorem and persistent homology of asymmetric
networks}, Journal of Applied and Computational Topology 2 (2018), 115--175.
\url{https://doi.org/10.1007/s41468-018-0020-6}.

\bibitem{BrunSalbu}
M. Brun and L. M. Salbu, \emph{The Rectangle Complex of a Relation},
Mediterranean Journal of Mathematics 20 (2023), article 7,
Theorems 4.3, 5.1, and 5.2.
\url{https://doi.org/10.1007/s00009-022-02213-0}.

\bibitem{BauerEtAl}
U. Bauer, M. Kerber, F. Roll, and A. Rolle,
\emph{A Unified View on the Functorial Nerve Theorem and its Variations},
Expositiones Mathematicae 41(4) (2023), article 125503;
updated preprint v6, 2025.
\url{https://arxiv.org/abs/2203.03571}.

\bibitem{Hatcher}
A. Hatcher, \emph{Algebraic Topology}, Cambridge University Press, 2002.
\url{https://pi.math.cornell.edu/~hatcher/AT/AT.pdf}.

\bibitem{Erdmann}
M. Erdmann, \emph{Topology of Privacy: Lattice Structures and Information
Bubbles for Inference and Obfuscation}, 2017.
\url{https://www.cs.cmu.edu/~me/papers/TopPrivacy.pdf}.
\url{https://arxiv.org/abs/1712.04130}.

\bibitem{GreenEtAl}
R. E. Green, C. A. Joslyn, A. Myers, M. G. Rawson, and M. Robinson,
\emph{The Topological Structures of the Orders of Hypergraphs},
arXiv:2504.11760v2, 18 April 2025.
\url{https://arxiv.org/html/2504.11760v2}.

\bibitem{Stone}
M. H. Stone, \emph{The Theory of Representations for Boolean Algebras},
Transactions of the American Mathematical Society 40 (1936), 37--111.
\url{https://www.ams.org/tran/1936-040-01/S0002-9947-1936-1501865-8/S0002-9947-1936-1501865-8.pdf}.

\bibitem{HermannKrener}
R. Hermann and A. J. Krener,
\emph{Nonlinear Controllability and Observability},
IEEE Transactions on Automatic Control 22(5) (1977), 728--740.
\url{https://www.math.ucdavis.edu/~krener/1-25/10.IEEETAC77.pdf}.

\bibitem{ShaliziCrutchfield}
C. R. Shalizi and J. P. Crutchfield,
\emph{Computational Mechanics: Pattern and Prediction, Structure and
Simplicity}, Journal of Statistical Physics 104 (2001), 817--879.
\url{https://arxiv.org/abs/cond-mat/9907176}.

\bibitem{Xue}
T. Xue, \emph{Definitional Extension in Type Theory},
TYPES 2013, LIPIcs 26 (2014), 251--269, Section 3.
\url{https://doi.org/10.4230/LIPIcs.TYPES.2013.251}.

\bibitem{GengelbachWeber}
A. Gengelbach and T. Weber,
\emph{Model-Theoretic Conservative Extension for Definitional Theories},
LSFA 2017, Electronic Notes in Theoretical Computer Science 338, 133--145.
\url{https://doi.org/10.1016/j.entcs.2018.10.009}.

\bibitem{Lawvere}
F. W. Lawvere, \emph{Functorial Semantics of Algebraic Theories},
1963, reprinted in Reprints in Theory and Applications of Categories 5 (2004).
\url{https://www.tac.mta.ca/tac/reprints/articles/5/tr5.pdf}.

\bibitem{GoguenBurstall}
J. A. Goguen and R. M. Burstall,
\emph{Institutions: Abstract Model Theory for Specification and Programming},
Journal of the ACM 39(1) (1992), 95--146.
\url{https://doi.org/10.1145/147508.147524}.

\bibitem{Caramello}
O. Caramello, \emph{The unification of mathematics via topos theory},
2010; author materials on classifying toposes and bridges.
\url{https://arxiv.org/abs/1006.3930}.
\url{https://www.oliviacaramello.com/Unification/UnilogTutorial.html}.
\end{thebibliography}

\end{document}
